\twocolumn[{
\section{Failure Errors by Phase and Reason}
\begin{center}\small
\begin{tabular}{llrrll}
\toprule
Phase & Reason group & $n$ & Obs./Pred. & Baseline errors & Current blend errors \\
\midrule
I & Safety & 10 & 1/9 & 4/10 (40.0\%) & 5/10 (50.0\%) \\
I & Funding & 4 & 1/3 & 2/4 (50.0\%) & 2/4 (50.0\%) \\
I & Efficacy & 5 & 0/5 & 5/5 (100.0\%) & 5/5 (100.0\%) \\
I & Recruitment & 22 & 2/20 & 10/22 (45.5\%) & 7/22 (31.8\%) \\
I & Other & 2 & 0/2 & 1/2 (50.0\%) & 1/2 (50.0\%) \\
I & No recorded reason & 28 & -- & 17/28 (60.7\%) & 14/28 (50.0\%) \\
\midrule
II & Safety & 12 & 1/11 & 5/12 (41.7\%) & 5/12 (41.7\%) \\
II & Funding & 4 & 1/3 & 0/4 (0.0\%) & 0/4 (0.0\%) \\
II & Efficacy & 18 & 0/18 & 4/18 (22.2\%) & 5/18 (27.8\%) \\
II & Recruitment & 91 & 12/79 & 17/91 (18.7\%) & 17/91 (18.7\%) \\
II & Other & 15 & 3/12 & 5/15 (33.3\%) & 4/15 (26.7\%) \\
II & No recorded reason & 144 & -- & 48/144 (33.3\%) & 48/144 (33.3\%) \\
\midrule
III & Safety & 10 & 1/9 & 7/10 (70.0\%) & 7/10 (70.0\%) \\
III & Funding & 0 & 0/0 & -- & -- \\
III & Efficacy & 10 & 0/10 & 10/10 (100.0\%) & 10/10 (100.0\%) \\
III & Recruitment & 23 & 5/18 & 14/23 (60.9\%) & 14/23 (60.9\%) \\
III & Other & 6 & 2/4 & 5/6 (83.3\%) & 5/6 (83.3\%) \\
III & No recorded reason & 82 & -- & 52/82 (63.4\%) & 50/82 (61.0\%) \\
\bottomrule
\end{tabular}

\captionof{table}{True outcome failures predicted successful at threshold 0.5: $k/n$ and percentage.
Both models use the same historical Part 1 reason lookup. ``Obs./Pred.'' counts observed versus
estimated reason memberships; groups overlap. Unmatched trials form the no-recorded-reason
group. A dash indicates no eligible trials or no applicable reason source.}
\label{tab:errors}
\end{center}
}]
\subsection{Definition and provenance}
For each phase/reason group $g$ containing only true failures,
\begin{equation}
 \mathrm{FSR}_{g}=\frac{\sum_{i\in g}\mathbb{I}[p_i\geq0.5]}{|g|}.
\end{equation}
A trial counts once within each phase/group, but may belong to several groups.
An absent Part 1 record is not a negative reason label.

Freezing the historical K-Means/LR reason lookup isolates changes in outcome predictions.
This is not an analysis of the newer 235-label ridge reason model. We distinguish observed,
estimated, and training-resubstitution reasons.

The 452 failed trial-phase rows represent 398 unique trials. Of 198 matched rows, 27 have observed reasons and 171 have predicted reasons; 254 rows have no Part 1 match. Among predicted rows, 131 are training resubstitution estimates. These are not out-of-fold or confirmed labels. Matched empty-reason groups have zero support after fallback and are omitted.


\subsection{Interpretation}
Phase I false-success counts change from 34/64 to 29/64, while true-success false negatives change from 9 to 13.

Phase II false-success counts change from 73/265 to 73/265, while true-success false negatives change from 96 to 94.

Phase III false-success counts change from 80/123 to 78/123, while true-success false negatives change from 21 to 20.

Phase I Recruitment errors fall from 10/22 to 7/22, while Safety errors rise from 4/10 to 5/10. Phase II Other improves from 5/15 to 4/15, but Efficacy worsens from 4/18 to 5/18. All five Phase I and all ten Phase III Efficacy failures remain false successes. These groups include estimated reasons.


Small groups make one-trial changes look large; Phase III has no Funding failures.
Estimated reasons add uncertainty beyond sampling variability. These descriptive associations
do not establish causes of errors. Stopping reasons never enter the outcome model as features.
